\documentclass[10pt]{article}
\usepackage[T1]{fontenc}
\usepackage[a4paper,margin=19mm]{geometry}
\usepackage{amsmath,amssymb,booktabs,array}
\usepackage{enumitem}
\usepackage{xspace}
\usepackage{xcolor}
\usepackage[colorlinks=true,linkcolor=blue!45!black,urlcolor=blue!45!black]{hyperref}
\definecolor{accent}{HTML}{315B7D}
\setlength{\parindent}{0pt}
\setlength{\parskip}{4pt}
\setlist{nosep,leftmargin=*}
\setcounter{secnumdepth}{2}
\newcommand{\goal}[1]{\medskip\noindent\textcolor{accent}{\textbf{Experimental goal.}} #1\medskip}
\newcommand{\code}[1]{\texttt{#1}}
\newcommand{\Tzero}{\ensuremath{T_0}\xspace}
\newcommand{\Tone}{\ensuremath{T_1}\xspace}
\newcommand{\Ttwo}{\ensuremath{T_2}\xspace}
\newcommand{\Tthree}{\ensuremath{T_3}\xspace}
\pagestyle{plain}

\title{Adaptation Experiment 3: Learned Retrieval Gates}
\author{}
\date{}
\begin{document}
\maketitle

\section{Candidates and targets}
Two gate families are evaluated. A cov gate chooses between the vanilla
forecast $V$ and the covariate-informed forecast $C$. An avgy
gate chooses between $V$ and
\[
 \overline Y_h=\sum_{j=1}^{k}w_jY_{j,h}.
\]
Write either candidate as $A\in\{C,\overline Y\}$. The horizon-level advantage
used for supervision is
\[
 \Delta_h=(Y_h-V_h)^2-(Y_h-A_h)^2.
\]
A positive value means the candidate is better. Shared gates make one decision
for the complete forecast window using $\overline\Delta=H^{-1}\sum_h\Delta_h$;
horizon gates make $H$ separate decisions.

\goal{Learn when cov or avgy should replace the reliable
vanilla forecast, and compare conditional gates with no-feature and
target-aware upper bounds.}

\section{Engineered features}
Let
\[
 \mu_w(z)=\sum_jw_jz_j,\qquad
 \sigma_w(z)=\sqrt{\sum_jw_j(z_j-\mu_w(z))^2}.
\]
For each neighbor define the across-horizon summaries
\[
\begin{aligned}
 D^Y_j&=H^{-1}\sum_h(Y_{j,h}-V_h),
 &S^Y_j&=\operatorname{sd}_h(Y_{j,h}-V_h),\\
 D^N_j&=H^{-1}\sum_h(Y_{j,h}-N_{j,h}),
 &S^N_j&=\operatorname{sd}_h(Y_{j,h}-N_{j,h}).
\end{aligned}
\]
Both gate shapes receive
\[
 \mu_w(D^Y),\ \sigma_w(D^Y),\ \mu_w(S^Y),\qquad
 \mu_w(D^N),\ \sigma_w(D^N),\ \mu_w(S^N).
\]
They also receive the query mean and standard deviation; the mean and
between-neighbor standard deviation of neighbor look-back means and
look-back standard deviations; same-user ratio; mean and standard deviation of
neighbor age; weight standard deviation and maximum; and mean retrieval
distance.

The candidate features differ by gate shape:
\begin{itemize}
\item A shared gate receives
$\operatorname{mean}_h(A_h-V_h)$ and
$\operatorname{sd}_h(A_h-V_h)$.
\item The model for horizon $h$ receives only $A_h-V_h$, not the complete
vector $(A_1-V_1,\ldots,A_H-V_H)$. It additionally receives the local
$\mu_w(Y_{j,h}-V_h)$, $\sigma_w(Y_{j,h}-V_h)$,
$\mu_w(Y_{j,h}-N_{j,h})$, and $\sigma_w(Y_{j,h}-N_{j,h})$.
\end{itemize}
Thus every horizon gate has the same feature count and there are $H$
independent models; information from another candidate horizon cannot leak
through the candidate-difference feature.

\section{Models and artifact names}
\small
\begin{tabularx}{\linewidth}{p{.34\linewidth}X}
\toprule
Name pattern & Meaning \\
\midrule
\code{bayes\_<A>\_shared/horizon}
 & No-feature decision from the average \Tone+\Ttwo advantage.\\
\code{catboost\_<A>\_classifier\_}\\[-.25ex]
\code{shared/horizon}
 & CatBoost estimate of $\mathbb 1[\Delta>0]$; choose $A$ above probability
 $1/2$.\\
\code{catboost\_<A>\_regressor\_}\\[-.25ex]
\code{shared/horizon}
 & CatBoost estimate of the signed loss advantage; choose $A$ above zero.\\
\code{oracle\_<A>\_shared/horizon}
 & Target-aware upper bound evaluated on the scored split; never deployable.\\
\bottomrule
\end{tabularx}
\normalsize
Here \code{<A>} is \code{cov} or \code{avgy}.

The primary screen retains only
\code{catboost\_<A>\_regressor\_shared}. Its signed-advantage target matches
the metric being optimized and one decision is made for the complete forecast
window. The shared no-feature Bayes decision and shared oracle remain in that
screen as lower-complexity and target-aware references.
Its primary screen dataset set is Electricity, Traffic, Solar, and Exchange.
The final full/ultra benchmark transfers only selected complete gate pipelines
to those datasets plus the four quantity-separated ETT temperature/load
panels, while retaining the three screen settings. Original ETTh1, ETTh2, ETTm1,
ETTm2 and Weather are evaluated only by
\code{mixed\_quantity\_ablation.slurm}, since their channels mix unlike
physical quantities and can make cross-channel retrieval semantically
ambiguous.

The root front \code{catboost\_ablation.slurm} is a separate later study. It
accepts selected shared-regressor screen winners, holds each winner's candidate
and retrieval pipeline fixed, and evaluates all four CatBoost
formulations---classifier/regressor crossed with shared/per-horizon decisions.
It also reruns matching shared/per-horizon Bayes and oracle references, so the
ablation is self-contained without reopening the full retrieval screen.

\section{Validation protocol}
Extraction provides pooled \Tone+\Ttwo and fixed \Tthree payloads. For these
gates, the last 20\% of pooled query dates is \Ttwo by default:
\begin{enumerate}
\item fit CatBoost on \Tone with at most 300 trees, learning rate 0.03,
depth 4, and balanced classifier classes;
\item use \Ttwo only for early stopping and selection of the tree count;
\item instantiate a fresh model with that selected count and refit on all
of \Tone+\Ttwo;
\item score once on complete \Tthree.
\end{enumerate}
CatBoost does not require a separate hyperparameter search here: depth and
learning rate are fixed experimental choices, while \Ttwo selects effective
capacity through early stopping. Feature importance is saved separately for
each candidate, objective, and decision shape.

The implementation fits horizons serially.  For every gate it selects the tree
count, releases the selection model, refits one final model, writes that model
in CatBoost's native format, scores disk-backed arrays, and releases the model
before continuing.  Gate summaries are computed in bounded row chunks and a
single horizon feature matrix is materialized at a time.  Predictions and
diagnostic scores are separate NumPy arrays indexed by the sole current
\code{prediction\_manifest.json}; \code{gate\_artifacts.json} indexes the
models, and \code{result\_manifest.json} is written last as the completion
marker.

This refit is valid because $C$ and $\widehat Y$ are fixed candidates. For a
future gate whose candidate is itself trained, the stage-two protocol is
different: fit the candidate on \Tone and the gate on \Ttwo without refitting
the candidate afterward. Refitting the candidate on \Tone+\Ttwo would change
the gate inputs and make its \Ttwo supervision in-sample. Avoiding that bias
requires out-of-fold candidate predictions or an additional split.

\textbf{Launchers:} \code{gates.slurm}; \code{catboost\_ablation.slurm}\\
\textbf{Module:} \code{src.adaptors.baselines.evaluate --family gates}

\section{Interpretation}
The no-feature gate measures whether one candidate wins globally. CatBoost
adds conditional decisions. The oracle gap is headroom, not performance.
Context and aggregated-neighbor gates answer different questions and should
both be reported against $V$ on identical \Tthree windows.
\end{document}
